

\section{Probabilistic Boolean tapes}\label{sec:probbooltapes}
In \cite[Example 30]{bonchi2025tapediagramsmonoidalmonads}, an encoding of probabilistic Boolean circuits into tape diagrams is presented. 
In this section, we first recall such encoding (Section~\ref{ssec:probbooltapes:encoding}), we then illustrate an axiomatisation (Section~\ref{ssec:probbooltapes:axiomatisation}) and we prove its completeness (Section~\ref{ssec:probbooltapes:completeness}). 



\subsection{Encoding probabilistic Boolean circuits into tapes}\label{ssec:probbooltapes:encoding}
Example 30 in \cite{bonchi2025tapediagramsmonoidalmonads} encodes probabilistic Boolean circuits from $\Diag{\SigPRB}$ into $\CatT{\Diag{\SigPB}}$, namely probabilistic tape diagrams of \emph{partial} Boolean circuits.
%
The encoding is defined as the unique (symmetric strict) monoidal functor $\encoding{-} \colon (\Diag{\SigPRB},\otimes, \uno) \to (\CatT{\Diag{\SigPB}},\otimes, \uno)$ mapping
\[
\resizebox{0.95\linewidth}{!}{$
\Andgate \mapsto \Andgate[t] \qquad
\Notgate \mapsto \Notgate[t] \qquad
\Flip{1} \mapsto \Flip{1}[t] \qquad
\CBcopier \mapsto  \CBcopier[t] \qquad
\CBdischarger \mapsto \CBdischarger[t] \qquad
\CBcocopier \mapsto \CBcocopier[t] 
$}
\]
and $\Flip{p} \mapsto \Flip{p}[t]$ where 
\begin{equation}\label{eq:exprobabilistictapes}
\Flip{p}[t] \defeq \tikzfig{tapes/examples/1p0}\end{equation}

The abstract theory developed in \cite{bonchi2025tapediagramsmonoidalmonads} guarantees the existence of a unique morphism of rig categories $\dsem{-}\colon \CatT{\Diag{\SigPB}} \to \KlD$ which assigns to each tape its semantics in $\KlD$. For convenience of the reader the definition of $\dsem{-}$  is reported in Table~\ref{eq:SEMANTICA}. For example, $\dsem{-}$ assigns to  the following tape 
\begin{equation*}
        \tikzfig{tapes/examples/ANDpOR} 
% \text{INSERT TAPE }1+_p 0\text{ ---- INSERT TAPE }AND +_p OR 
\end{equation*}
the arrow $2 \times 2 \to 2$ of \(\KlD\) which maps any pair of Booleans $(x,y)$ into $x\wedge y$ with probability $p$ and $x \vee y$ with probability $(1-p)$.
Instead, $\Flip{p}[t]$  is mapped by $\dsem{-}$ into the arrow $1 \to 2$ assigning to $\bullet$ the distribution $1 \mapsto p$, $0 \mapsto (1-p)$. Observe that this is exactly $\osem{\Flip{p}}$ as defined in \eqref{eq:semanticsFlip}. More generally, a simple inductive argument confirms that $\dsem{\encoding{c}}=\osem{c}$ for all $c$ in $\Diag{ PB}$, i.e., the encoding preserves the semantics.

\begin{proposition}\label{prop:encoding} For all $c\in \Diag{\SigPRB}$, $\osem{c}=\CBdsem{\encoding{c}}$.
\end{proposition}
\medskip

Interestingly, tapes in $\CatT{\Diag{\SigPB}}$ are strictly more expressive than probabilistic Boolean circuits in $\Diag{\SigPRB}$, in the sense that they denote a strictly larger class of arrows in $\KlD$. This difference becomes apparent by recalling, from~\cite{bonchi2025tapediagramsmonoidalmonads}, how probabilistic control is implemented in $\Diag{\SigPRB}$ and in $\CatT{\Diag{\SigPB}}$.

In~\cite{piedeleu2025boolean}, probabilistic control is realised via a multiplexer, represented by the tape $\scalebox{0.7}{\Ifgate[t]}$. Intuitively, when $\Flip{p}[t]$, $\Tcirc{c}{\!\!\!\!}{\!\!\!\!}$ and $\Tcirc{d}{\!\!\!\!}{\!\!\!\!}$ are connected, respectively, to the first, second, and third inputs of the multiplexer, the resulting output coincides with that of $c$ with probability $p$ and with that of $d$ with probability $1-p$.
Formally, this behaviour is captured by the composite
\[
(\Flip{p}[t] \;\per\; \Tcirc{c}{\!\!\!\!}{\!\!\!\!} \;\per\; \Tcirc{d}{\!\!\!\!}{\!\!\!\!}) \; ; \; \scalebox{0.7}{\Ifgate[t]},
\]
which, by the definition of $\Flip{p}[t]$ in~\eqref{eq:exprobabilistictapes} and of $\per$ in Table~\ref{tab:producttape}, corresponds to the tape shown on the left below:
\begin{equation*}
    \tikzfig{tapes/examples/multiplexT}
    \qquad \qquad
    \tikzfig{tapes/examples/pchoice}
\end{equation*}

Although the two diagrams above exhibit similar behaviour, a crucial difference emerges when $d$ (or symmetrically $c$) is instantiated as $\Flip{\bot}[t]$, which denotes the null subdistribution and intuitively represents failure. In this case, the circuit on the left always fails (Lemma~4.4 in~\cite{piedeleu2025boolean}). By contrast, the circuit on the right still produces the output of $c$ (respectively $d$) with probability $p$ (respectively $1-p$).

\begin{table}[t]
    \renewcommand{\arraystretch}{1.5}
\!\!\!\begin{tabular}{l@{\;\;\;\;}l@{\;\;\;\;}l@{\;\;\;\;}l@{\;\;\;\;}l}
$\CBdsem{\id{A}} \defeq \id{2} $&$ \CBdsem{\id{1}} \defeq \id{1} $&$ \CBdsem{\symmt{A}{A}}  \defeq \symmt{2}{2}  $&$ \CBdsem{c;d}  \defeq \CBdsem{c} ; \CBdsem{d}   $ & $ \CBdsem{c\per d}  \defeq \CBdsem{c}  \per \CBdsem{d} $\\
$\CBdsem{s}  \defeq \osem{s} $& $\CBdsem{\, \tapeFunct{c} \,} \defeq \CBdsem{c}  $ & $\CBdsem{\diagp{A^n}}  \defeq \; \diagp{2^n}$& $\CBdsem{\bang{A^n}} \defeq \bang{2^n}$  \; $\CBdsem{\cobang{A^n}}  \defeq \cobang{2^n}$    & $\CBdsem{\codiag{A^n}} \defeq \codiag{2^n} $ \\
$\CBdsem{\id{A^n}} \defeq \id{2^n} $&$ \CBdsem{\id{\zero}}  \defeq \id{\zero}  $&$ \CBdsem{\symmp{A^n}{A^m}} \defeq \symmp{2^n}{2^m} $&$ \CBdsem{\s;\t}  \defeq \CBdsem{\s}  ; \CBdsem{\t}  $&$ \CBdsem{\s \piu \t}  \defeq \CBdsem{\s}  \piu \CBdsem{\t}  $ % \\
% & & $\CBdsem{\bang{A^n}} \defeq \bang{2^n}$  
% \multicolumn{5}{c}{
%     \begin{tabular}{c@{\qquad\qquad}c}
%         $\CBdsem{\copier{A}} = \copier{\alpha_\sort(A)}$ & $\CBdsem{\discharger{A}} = \discharger{\alpha_\sort(A)}$
%     \end{tabular}
% }
\end{tabular}
\caption{The semantics $\dsem{-}\colon \CatT{\Diag{\SigPB}} \to \KlD$. Here $s$ is a generator in $\SigPB$.}\label{eq:SEMANTICA}
\end{table}

\medskip





\subsection{Axiomatisation}\label{ssec:probbooltapes:axiomatisation}
In order to provide a complete axiomatisation for the equivalence induced by $\dsem{-}$, we reuse $\eqPB$ which, as proved in Theorem \ref{thm:completenesspartialcircuits}, provides a complete axiomatisation for partial Boolean circuits. 
Since the axioms in $\ThPB$ are sound with respect to $\osem{-}$ and since $\osem{-}=\dsem{\encoding{-}}$, these axioms are also sound for $\dsem{-}$. However, they are not complete, as explained below.

\begin{example}
Consider the congruence $\tapeeqPB$ on $ \CatT{\Diag{\SigPB}}$ generated by $\eqPB$ as in \eqref{eq:congr2}. By Proposition \ref{prop:quotienttheory}, the quotient of $ \CatT{\Diag{\SigPB}}$ by  $\tapeeqPB$ is exactly $ \CatT{\Diag{\ThPB}}$ which, 
by Corollary~\ref{cor:isotapematrices} and Proposition~\ref{prop:isopartialboolean} is isomorphic to $\stmat{\Cat{Par}_2^+}$. Now consider only $1\times 1$ matrices: these are arrows in $\Cat{Par}_2^+$, i.e., subdistributions of partial functions. These are not in one-to-one correspondence with Kleisli arrows. Take for instance the null subdistribution $\star \in \Dis(\Cat{Par}[1, 2])$ and $\delta_\bot \in \Dis(\Cat{Par}[1, 2])$ which are distinct arrows in $\Cat{Par}_2^+[1,2]$. The former is drawn as the tape \scalebox{0.6}{\tikzfig{tapes/cipriano/AXPBP2right}} and the latter as $\Flip{\bot}[t]$. It is easy to see that $\dsem{-}$ maps both diagrams into the $\KlD$-arrow $\star_{1,2}$.
\end{example}



To obtain a complete axiomatisation, we need to add to $\ThPB$ the axioms in Figure~\ref{ax:BooleanTAPES}. The tape diagram on the right-hand side of \eqref{ax:AXPBP1} with probability $\frac{1}{2}$ tests if the input is $1$ and returns $1$ and, with probability $\frac{1}{2}$, tests if the input is $0$ and returns $0$. The equality \eqref{ax:AXPBP1} forces such tape to be equal to $\frac{1}{2} \cdot \id{A}$; Axiom \eqref{ax:AXPBP2} forces $\Flip{\bot}[t]$ to be the null subdistribution; Axiom \eqref{ax:canc} is an implication which asserts cancellativity: see Section~\ref{ssec:pca}. %Since cancellativity holds in $\KlD$, the reader can easily check that the axioms are sounds by computing the semantics of the left and right hand sides of the axioms \eqref{ax:AXPBP1} and \eqref{ax:AXPBP2}.

\begin{figure}[H]
  \centering

  % Labels
  \mylabelbis{ax:AXPBP1}{T2}
  \mylabelbis{ax:AXPBP2}{T1}
  \mylabelbis{ax:canc}{T3}
  
  \[
  \begin{array}{c@{\qquad}c@{\qquad}c}
    % Prima riga (due assiomi)
    
    \tikzfig{tapes/cipriano/AXPBP2left}
    \stackrel{\eqref{ax:AXPBP2}}{=}
    \tikzfig{tapes/cipriano/AXPBP2right}
    &
    \tikzfig{tapes/cipriano/AXPBP1left}
    \stackrel{\eqref{ax:AXPBP1}}{=}
    \tikzfig{tapes/cipriano/AXPBP1right}
    &
    {}
    \\[30pt]
    % Seconda riga (canc centrato)
    \multicolumn{3}{c}{
      \tikzfig{tapes/cipriano/Axcancleft}
      \stackrel{\eqref{ax:canc}}{\implies}
      \tikzfig{tapes/cipriano/Axcancright}
    }
  \end{array}
  \]
  
\caption{Axioms for tapes of partial Boolean circuits.}\label{ax:BooleanTAPES}
\end{figure}


From $\ThPB$ and the axioms in Figure~\ref{ax:BooleanTAPES}, one generates the congruence $\dot{\sim}$ on $\CatT{\Diag{\SigPB}}$ by means of the following inference rules.
\begin{equation}\label{eq:congr3}
        \scalebox{0.85}{$
        \centering
        \begin{array}{c}
        \begin{array}{c@{\qquad\qquad}c@{\qquad\qquad}c@{\qquad\qquad}c@{\qquad\qquad}c}
            \inferrule*[right=\eqref{ax:AXPBP1}]{\t_1 \stackrel{\eqref{ax:AXPBP1}}{=} \t_2}{\t_1 \dot{\sim} \t_2}
            &
            \inferrule*[right=\eqref{ax:AXPBP2}]{\t_1 \stackrel{\eqref{ax:AXPBP2}}{=} \t_2}{\t_1 \dot{\sim} \t_2}
            &
            \inferrule*[right=($\textsc{R}$)]{-}{\t \dot{\sim} \t}
            &
            \inferrule*[right=($\textsc{S}$)]{\t_1 \dot{\sim} \t_2}{\t_2 \dot{\sim} \t_1}
            &    
            \inferrule*[right=($\textsc{T}$)]{\t_1 \dot{\sim} \t_2 \quad \t_2 \dot{\sim} \t_3}{\t_1 \dot{\sim} \t_3}
            % &
            % \inferrule*[right=($E$)]{ (t,t') \in E \and U \in \sort^\star}{ t_U \dot{\sim} t'_U}
        \end{array}
        \\[10pt]
        \begin{array}{c@{\qquad}c@{\qquad}c@{\qquad}c@{\qquad}c}
        \inferrule*[right=($\ThPB$)]{c \eqPB d}{\tape{c} \dot{\sim} \tape{d}}
            &
            \inferrule*[right=\eqref{ax:canc}]{p\cdot \t \dot{\sim} p\cdot \s }{\t \dot{\sim} \s}
            &
            \inferrule*[right=($;$)]{\t_1 \dot{\sim} \t_2 \quad \s_1 \dot{\sim} \s_2}{\t_1;\s_1 \dot{\sim} \t_2;\s_2}
            &
            \inferrule*[right=($\piu$)]{\t_1 \dot{\sim} \t_2 \quad \s_1 \dot{\sim} \s_2}{\t_1\piu\s_1 \dot{\sim} \t_2 \piu \s_2}
            &
            \inferrule*[right=($\per $)]{\t_1 \dot{\sim} \t_2 \quad \s_1 \dot{\sim} \s_2}{\t_1\per \s_1 \dot{\sim} \t_2 \per \s_2}       
        \end{array}
        \end{array}
        $}
\end{equation}

Simple computations confirm that the axioms are sound.
\begin{proposition}[Soundness]\label{prop:soundnesstapes}
For all $\s,\t \in \CatT{\Diag{\SigPB}}$, if $\s \eqsynbis \t$ then $\dsem{\s}=\dsem{\t}$.
\end{proposition}


To prove completeness, it is convenient to consider diagrams in $\CatT{\Diag{\ThPB}}$. Let $\eqsyn$ be the congruence on $\CatT{\Diag{\ThPB}}$ defined as $\dot{\sim}$ but omitting the rule ($\ThPB$). By means of Proposition~\ref{prop:quotienttheory}, one can easily prove the following result.



\begin{proposition}\label{prop:twoequivalence}
For all $\s , \t \in \CatT{\Diag{\SigPB}}$, $\s \dot{\sim} \t$ iff $\CatT{Q_{\ThPB}}(\s) \eqsyn \CatT{Q_{\ThPB}}(\t)$.
\end{proposition}




Thanks to the above proposition, we can work with $\eqsyn$ rather than $\dot{\sim}$. We conclude this section by illustrating three key properties of $\eqsyn$. First, axiom \eqref{ax:AXPBP2} easily entails the following.

\begin{lemma}\label{lemma:bottomstar2}
%\label{axpbp2m}
 $\star_{A^0,A^n} \sim \Flipn{\bot}[t]$.  
\end{lemma}


Hereafter, we write $\Cat{B}[1,A^n]$ for the set of $n$-ary Boolean vectors, i.e., all those tapes in $\CatT{\Diag{\ThPB}}[1,A^n]$ of the form $\overrightarrow{b}\defeq\bigotimes_{i=1}^n b_i$ for $b_i \in \{\Flip{0}[t], \Flip{1}[t]\}$ if $n\not=0$, and $\tikzfig{tapes/empty}$ if $n=0$. If $\overrightarrow{b}\in\Cat{B}[1,A^n]$, we write $\overleftarrow{b}$ for the tape $\bigotimes_{i=1}^n b'_i\in \CatT{\Diag{\ThPB}}[A^n,1]$ where $b'_i$ is $\coflip{1}[t]$ if $b_i$ is $\Flip{1}[t]$ and $\coflip{0}[t]$ if $b_i$ is $\Flip{0}[t]$.

The second property crucially relies on axiom \eqref{ax:AXPBP1}.

\begin{lemma}\label{lemma:axiomsninput}
%\label{axpbp1m}
 For all $n\in \mathbb{N}$, $\frac{1}{2^n}\cdot \id{A^n} \sim \sum_{\overrightarrow{b}\in \Cat{B}[1,A^n]}^{} \frac{1}{2^n}\cdot\overleftarrow{b};\overrightarrow{b}$;
\end{lemma}
\begin{proof}
By induction  on $n$. The base case is trivial. For the inductive step, we have
 \begin{align}
  \frac{1}{2^{n+1}}\cdot \id{A^{n+1}}&= \frac{1}{2^n}\frac{1}{2}\cdot( \id{A}\otimes \id{A^n}) \notag\\
  &= \frac{1}{2^n}\cdot\left(\frac{1}{2}\cdot \id{A}\otimes \id{A^n} \right) \tag{\eqref{eq:monoidalenrichment}}\\
  &= (\frac{1}{2}\cdot \id{A})\otimes (\frac{1}{2^n}\cdot \id{A^n}) \tag{\eqref{eq:monoidalenrichment}}\\
  &\sim (\frac{1}{2}\cdot \id{A})\otimes \left(\sum_{\overrightarrow{b}\in \Cat{B}[1,A^n]}^{} \frac{1}{2^n}\cdot\overleftarrow{b};\overrightarrow{b}\right) \tag{Inductive hypothesis}\\
  &\sim (\coflip{1}[t];\Flip{1}[t] +_{\frac{1}{2}} \coflip{0}[t];\Flip{0}[t]) \otimes \left(\sum_{\overrightarrow{b}\in \Cat{B}[1,A^n]}^{} \frac{1}{2^n}\cdot\overleftarrow{b};\overrightarrow{b}\right) \tag{Axiom \ref{ax:AXPBP1}}\\
  &= \sum_{\overrightarrow{b}\in \Cat{B}[1,A^n]}^{} \frac{1}{2^n}\cdot\left((\coflip{1}[t];\Flip{1}[t] +_{\frac{1}{2}} \coflip{0}[t];\Flip{0}[t]) \otimes (\overleftarrow{b};\overrightarrow{b})\right) \tag{\eqref{eq:monoidalenrichment}}\\
  &=\sum_{\overrightarrow{b}\in \Cat{B}[1,A^n]}^{} \frac{1}{2^n}\cdot\left(((\coflip{1}[t];\Flip{1}[t]) \otimes (\overleftarrow{b};\overrightarrow{b})) +_{\frac{1}{2}} ((\coflip{0}[t];\Flip{0}[t]) \otimes (\overleftarrow{b};\overrightarrow{b}))\right) \tag{SMC}\\
  &= \sum_{\overrightarrow{b}\in \Cat{B}[1,A^n]}^{} \frac{1}{2^n}\cdot\left((\coflip{1}[t]\otimes \overleftarrow{b});(\Flip{1}[t]\otimes \overrightarrow{b}) +_{\frac{1}{2}} (\coflip{0}[t]\otimes \overleftarrow{b});(\Flip{0}[t]\otimes \overrightarrow{b})\right) \tag{\eqref{eq:monoidalenrichment}}
 \end{align}
 Now, since the set $\{\Flip{1}[t]\otimes \overrightarrow{b}, \Flip{0}[t]\otimes \overrightarrow{b}\}$ as $\overrightarrow{b}\in \Cat{B}[1,A^n]$ is exactly $\Cat{B}[1,A^{n+1}]$, then, rearranging the above sum we obtain $\frac{1}{2^{n+1}}\cdot \id{A^{n+1}}= \sum_{\overrightarrow{b}\in \Cat{B}[1,A^{n+1}]}^{} \frac{1}{2^{n+1}}\cdot\overleftarrow{b};\overrightarrow{b}$, as desired.
\end{proof}

The  lemma above and axiom \eqref{ax:canc} give us the third key result.

\begin{lemma}\label{lemma:PBP2} 
Let $\s,\t\in \CatT{\Diag{\ThPB}}[A^n,A^m]$. If, for all $\overrightarrow{b}\in \Cat{B}[1,A^n]$, $\overrightarrow{b};\s \eqsyn\overrightarrow{b}; \t$, then $\s \eqsyn \t$.
\end{lemma}
\begin{proof}
Consider the following derivation: 
\begin{align}
  \frac{1}{2^n} \cdot \t &= \frac{1}{2^n}\cdot (\id{A^n};\t) \notag\\
  &= (\frac{1}{2^n}\cdot \id{A^n});\t \tag{Lemma \ref{lemma:initialproperties2}}\\ 
  &\sim (\sum_{\overrightarrow{b}\in \Cat{B}[1,A^n]}^{} \frac{1}{2^n}\cdot\overleftarrow{b};\overrightarrow{b});\t \tag{Lemma \ref{lemma:axiomsninput}}\\
  &= \sum_{\overrightarrow{b}\in \Cat{B}[1,A^n]}^{} \frac{1}{2^n}\cdot\overleftarrow{b};\overrightarrow{b};\t \tag{\eqref{eq:enr}}\\
  &\sim \sum_{\overrightarrow{b}\in \Cat{B}[1,A^n]}^{} \frac{1}{2^n}\cdot\overleftarrow{b};\overrightarrow{b};\s \tag{Hp.}\\
  &=(\sum_{\overrightarrow{b}\in \Cat{B}[1,A^n]}^{} \frac{1}{2^n}\cdot\overleftarrow{b};\overrightarrow{b});\s \tag{\eqref{eq:enr}}\\
  &\sim (\frac{1}{2^n}\cdot \id{A^n});\s \tag{Lemma \ref{lemma:axiomsninput}}\\
  &= \frac{1}{2^n} \cdot \s. \tag{Lemma \ref{lemma:initialproperties2}} 
\end{align}
Hence, applying axiom \eqref{ax:canc}, we conclude that $\t \eqsyn \s$.
\end{proof}



\subsection{Completeness}\label{ssec:probbooltapes:completeness}
We write $\eqsynq{\CatT{\Diag{\ThPB}}}$ for the category obtained as the quotient of $\CatT{\Diag{\ThPB}}$ by $\eqsyn$ and by $Q_\eqsyn\colon \CatT{\Diag{\ThPB}} \to \eqsynq{\CatT{\Diag{\ThPB}}}$ the functor mapping each tape into its $\eqsyn$ equivalence class. By Propositions~\ref{prop:soundnesstapes} and \ref{prop:twoequivalence}, $\dsem{-}\colon \CatT{\Diag{\SigPB}} \to \KlD$ can be factored as
\[\xymatrix{\CatT{\Diag{\SigPB}} \ar[r]^{\CatT{Q_\ThPB}} & \CatT{\Diag{\ThPB}} \ar[r]^{Q_\eqsyn} & \eqsynq{\CatT{\Diag{\ThPB}}} \ar@{.>}[r]^I& \KlD}\]
for some functor $I$. Proving completeness amounts to proving that  $I$  is faithful. %\colon \eqsynq{\CatT{\Diag{BPar}}} \to \KlD$

\medskip

%\marginpar{C:Quando diciamo che J=P;ParJ? inoltre non dovrebbe andare da Set?}
By Corollary~\ref{cor:isotapematrices} and Proposition~\ref{prop:isopartialboolean}  we know that 
\begin{equation}\label{eq:i2so}
  \CatT{\Diag{\ThPB}} \cong \stmat{\Diag{\ThPB}^+}\cong \stmat{\Cat{Par}_2^+}\text{.}
\end{equation}
%
Now, observe that the functor $\ParJ\colon \,\Cat{Par} \to \KlD$ from Section~\ref{ssec:probabilisticboolean} restricts to $\ParJ_2\colon \Cat{Par}_2 \to \KlD$. Since $\KlD$ is a convex biproduct category, by the two adjunctions in Theorems~ \ref{thm:freeenriched} and \ref{thm:matfree}, there exists a functor \[\ParJ_2'\colon \stmat{\Cat{Par}_2^+} \to \KlD\text{.} \]

%We call $F\colon \CatT{\Diag{\ThPB}} \to \stmat{\KlD_2}$, the morphism of convex biproduct categories obtained by composing the isomorphism $\CatT{\Diag{\ThPB}} \to  \stmat{\Cat{Par}_2^+}$ followed by $\stmat{\ParJ^\sharp}$. 
%\marginpar{C: prova di dimostrazione che fa a meno di $\J$.}

%We call $G\colon \stmat{\KlD_2} \to \KlD$ the morphism of convex biproduct categories obtained by composing the embedding $\stmat{\KlD_2} \to \stmat{\KlD}$ with the counit of the adjunction in Theorem~\ref{thm:matfree}, $\epsilon_{\KlD} \colon \stmat{\KlD}\to \KlD$. The latter is, by Proposition~\ref{prop:counit fullfaithful}, full and faithful and, therefore, so is $G$. 
%\[\xymatrix{
%\CatT{\Diag{\ThPB}}  \ar[r]^{\eqsynq{Q}} \ar[d]_{F} & \eqsynq{\CatT{\Diag{\ThPB}}} \ar[d]^{I} \ar@{.>}[dl]|H \\
%\stmat{\KlD_2} \ar[r]_G & \KlD
%}\]
%
\begin{comment}
We recall from \cite{arxivprobbooltapes} the following property of $\J^\sharp$.
\begin{lemma}\label{lemma:Jfaithful}
For all $d_1,d_2\in \Sets_2^+[1, 2^m]$, $\J^\sharp(d_1)=\J^\sharp(d_2)$ iff $d_1=d_2$.
\end{lemma}
\begin{proof}
It is convenient to give the explicit definition of $\J^\sharp$. Recall that an arrow $d\colon 2^n \to 2^m$ in $\Sets_2^+$ is a subprobability distribution on $\Sets_2[2^n,2^m]$, i.e., an element of $\Dis(\Sets_2[2^n,2^m])$. $\J^\sharp$ is the identity on objects functor mapping $d\colon 2^n \to 2^m$, into the Kleisli arrow defined for all $b\in 2^n, b'\in 2^m$ as $\J^\sharp(d)(b'|b)=\sum_{\{f\mid f(b)=b'\}}d(f)$. It is trivial to see that for $n=0$, $\J^\sharp$ provides the obvious bijection $$\Dis(\Sets_2[1,2^m])\cong \Dis(2^m) \cong   \KlD_2[1,2^m]\text{.}$$
%
%$b=\bullet$ (recall that $\bullet$ is the unique element of $1$), $J^\sharp(d)(b'|\bullet)=\sum_{\{f\mid f(\bullet)=b'\}}d(f)$.
%Define $U\colon \KlD_2[1,2^m] \to \Sets_2^+[1,2^m]$ as the function mapping any $g\colon 1 \to \Dis(2^m)$ into $\sum_{b'\in 2^m} g(b'|\bullet)\cdot \bar{b'}$ where $\bar{b'}\colon 1 \to 2^m$ is the function mapping $\bullet$ into $b'$. It is immediate to see that $U$ and $J^\sharp$ provides the obvious bijection $\Dis(\Sets_2[1,2^m])\cong \KlD_2[1,2^m]$.
\end{proof}
\end{comment}
%
%
%Since $\eqsynq{Q}$ is the identity on objects and $G$ is full and faithful, then by the properties of orthogonal factorisation systems (see e.g. \cite{nlab:bofactorizationsystem}), there exists a unique $H\colon  \eqsynq{\CatT{\Diag{\ThB}}} \to \stmat{\KlD_2}$ making the above diagram commute. Below we prove that $H$ is faithful to conclude that also $I$ is faithful.


%We first prove that $F$ is faithful on vectors of Booleans.
Consider the diagram below, where the vertical isomorphisms are those in \eqref{eq:i2so}.

\[
\qquad \xymatrix@R=2ex@C=3ex{
\CatT{\Diag{\ThPB}} \ar[r]^{\eqsynq{Q}} \ar[d]_{\cong} 
  & \eqsynq{\CatT{\Diag{\ThPB}}} \ar[dd]^{I} \\
{\stmat{\Diag{\ThPB}^+}} \ar[d]_{\cong} 
  & {} \\
\stmat{\Cat{Par}_2^+} \ar[r]_{\ParJ_2'} 
  & \KlD
}
\]
In the following result we prove that the above diagram commutes, where $\comp$ denotes the composition of the left vertical isomorphisms with $\ParJ_2'$.

\begin{lemma}\label{lemma:diagramma-commutativo}
$I\circ \eqsynq{Q}=\comp$.
\end{lemma}

\begin{lemma}\label{lemma:Ffaith}
For all $\t,\t'\in \CatT{\Diag{\ThB}}[A^0,A^m]$, if $\comp(\t)=\comp(\t')$, then $\t=\t'$.
\end{lemma}



Now, by relying on Lemma~\ref{lemma:caratterizzazionecircuiti0n}, we have the following normal forms for tapes $\t \colon A^0 \to A^m$.

\begin{lemma}\label{lemma:PBP0}
For all  $\t\in \CatT{\Diag{\ThPB}}[A^0,A^m]$, there exist $n\in\mathbb{N}$, $p_1,\dots,p_n\in (0,1)$, $\overrightarrow{b}_1,\dots, \overrightarrow{b}_n\in \Cat{B}[1,A^m]$ and $p,q\in [0,1]$, such that $\sum_{i=1}^{n}p_i= 1$, and $$\t=(\sum_{i=1}^{n}p_i\cdot \overrightarrow{b}_i)+_p(\star_{A^0,A^m} +_q \Flipm{\bot}[t])\text{.}$$ 
  \end{lemma}


The previous two results, combined with Lemma~\ref{lemma:bottomstar2} allow us to prove the following key property.



\begin{lemma}\label{lemma:PBP1}
  For all $\t,\t'\in \CatT{\Diag{\ThPB}}[A^0,A^m]$, if $\comp(\t)=\comp(\t')$, then $\t\sim \t'$.
\end{lemma}

\begin{theorem}\label{thm:completenessPBPtapes}
  $I\colon \eqsynq{\CatT{\Diag{\ThPB}}} \to \stmat{\KlD_2}$ is faithful.
\end{theorem}
\begin{proof}
  Let $[\t]_\sim, [\s]_\sim \in \eqsynq{\CatT{\Diag{\ThPB}}}[A^n,A^m]$ be such that $I([\t]_\sim)=I([\s]_\sim)$. Then: %\marginpar{Manca il vettore sopra i $b$?- C:sì}%, consider the following derivation:
  \begin{align}
    I([\t]_\sim)=I([\s]_\sim) &\implies \comp(\t)=\comp(\s) \tag{Lemma \ref{lemma:diagramma-commutativo}}\\
    &\implies \text{for all } \overrightarrow{b}\in \Cat{B}[1,A^n], \comp(\overrightarrow{b});\comp(\t)=\comp(\overrightarrow{b});\comp(\s) \notag \\
    &\implies \text{for all } \overrightarrow{b}\in \Cat{B}[1,A^n], \comp(\overrightarrow{b};\t)=\comp(\overrightarrow{b};\s) \tag{Functoriality}\\
    &\implies \text{for all } \overrightarrow{b}\in \Cat{B}[1,A^n], \overrightarrow{b};\t \eqsyn \overrightarrow{b};\s \tag{Lemma \ref{lemma:PBP1}}\\
    &\implies \t \eqsyn \s \tag{Lemma \ref{lemma:PBP2}}
  \end{align}
  Hence, $[\t]_\sim=[\s]_\sim$. The above derivation proves that $I$ is faithful for arrows of type $A^n \to A^m$.  This fact and convex products easily entail that $I$ is faithful for arrows $\s,\t \colon A^k \to \bigoplus_{i=1}^nA^{m_i}$. 
Indeed, by Lemma~\ref{decomposition} and Theorem~\ref{thm:TCconvexbiproductcategory}, 
\begin{equation}\label{eq:boh}\s=\langle \s_1, \dots, \s_n \rangle_{\vec{q}} \quad \t=\langle \t_1, \dots, \t_n \rangle_{\vec{p}}\end{equation} 
for some $\vec{q}=q_1,\dots, q_n$, $\vec{p}=p_1,\dots, p_n$, $\s_i \colon A^k \to A^{m_i}$ and $\t_i \colon A^k \to A^{m_i}$. Thus
\begin{align*}
I([\s]_\eqsyn) = I([\t]_\eqsyn) & \Rightarrow \comp(\s)=\comp(\t) \tag{Lemma \ref{lemma:diagramma-commutativo}}\\
& \Rightarrow \text{for all }i \text{, } \comp(\s);\comp(\pi_i)=\comp(\t);\comp(\pi_i) \\
& \Rightarrow \text{for all }i\text{, } \comp(\s;\pi_i)=\comp(\t;\pi_i) \tag{Functoriality} \\
& \Rightarrow  \text{for all }i\text{, }  \s;\pi_i\sim\t;\pi_i \tag{Previous implication}\\
& \Rightarrow  \text{for all }i\text{, }  q_i \cdot \s_i\sim p_i \cdot \t_i \tag{\ref{eq:boh}}\\
& \Rightarrow  \text{for all }i\text{, }  [q_i \cdot \s_i]_\sim = [p_i \cdot \t_i]_\sim 
\end{align*}
%To conclude that $[\s]_\sim=[\t]_\sim$, observe that 
Axiom \ref{ax:canc} allows us to apply Proposition~\ref{cor: quotient category is convex biproduct }, which states that $\CatT{\Diag{\ThPB}}_\sim$ is a convex biproduct category and that the quotient functor $Q_\eqsyn\colon \CatT{\Diag{\ThPB}} \to \CatT{\Diag{\ThPB}}_\sim$ is a morphism of convex biproduct categories.  Hence, Proposition~\ref{prop:functor 1 e mezzo} implies that $(\bigoplus_{i=1}^nA^{m_i}, [\pi_i]_\sim)$ is the $n$-ary convex product of $A^{m_1}, \dots, A^{m_n}$ in $\CatT{\Diag{\ThPB}}_\sim$ and $ [\s]_\sim$ is the unique arrow such that for all $i=1,\dots,n$,
\begin{align}
  [\s]_\sim;[\pi_i]_\sim &= [\s;\pi_i]_\sim \tag{Functoriality of $Q_\eqsyn$}\\
   &=[q_i \cdot \s_i]_\sim \tag{\ref{eq:boh}}\\
  &=q_i\cdot [\s_i]_\sim; \tag{pca-enrichment of $\CatT{\Diag{\ThPB}}_\sim$}
\end{align}
while $[\t]_\sim$ is the unique arrow such that for all $i=1,\dots,n$,
\begin{align}
  [\t]_\sim;[\pi_i]_\sim&= [\t;\pi_i]_\sim \tag{Functoriality of $Q_\eqsyn$}\\
  &= [p_i \cdot \t_i]_\sim \tag{\ref{eq:boh}}\\
  &=p_i\cdot [\t_i]_\sim. \tag{pca-enrichment of $\CatT{\Diag{\ThPB}}_\sim$}
\end{align}
Hence, since for all $i=1,\dots,n$, $[q_i \cdot \s_i]_\sim = [p_i \cdot \t_i]_\sim$, we have $[\s]_\sim = [\t]_\sim$.

%$[\s]_\sim=\langle [\s_1]_\sim, \dots, [\s_n]_\sim \rangle_{\vec{q}}$ and $[\t]_\sim=\langle [\t_1]_\sim, \dots, [\t_n]_\sim \rangle_{\vec{p}}$. Since for all $i$, $[\t];[\pi_i]_\sim$ and  $[\s];[\pi_i]_\sim$ are both equal to  $p_i\cdot[\t_i]_\sim=[p_i \cdot \t_i]_\sim=[q_i \cdot \s_i]_\sim$, the universal  property of  n-ary convex products in $\CatT{\Diag{\ThPB}}_\sim$ implies that $[\s]_\sim=[\t]_\sim $.

For arrows of arbitrary type $\bigoplus_{j=1}^o A^{k_j} \to \bigoplus_{i=1}^nA^{m_i}$, one can easily rely on the universal property of coproducts and the case that we just proved.
\end{proof}

\begin{corollary}\label{cor:completenessPBPtapes}
  For all $\s,\t \in \CatT{\Diag{\SigPB}}$, if $\dsem{\s}=\dsem{\t}$ then $\s \eqsynbis \t$.
\end{corollary}
\begin{proof}
By Theorem~\ref{thm:completenessPBPtapes}, $I$ is faithful. The following derivation concludes the proof.
  \begin{align}
    \dsem{\s}=\dsem{\t} &\Longleftrightarrow (\CatT{Q_{\ThPB}}; Q_\eqsyn ; I) (\s)= (\CatT{Q_{\ThPB}}; Q_\eqsyn ; I) (\t) \tag{$\dsem{-}= \CatT{Q_{\ThPB}}; Q_\eqsyn ; I$}\\
    &\implies   (\CatT{Q_{\ThPB}}; Q_\eqsyn ) (\s)= (\CatT{Q_{\ThPB}}; Q_\eqsyn )(\t) \tag{$I$ is faithful} \\
    &\Longleftrightarrow \CatT{Q_{\ThPB}}  (\s) \eqsyn \CatT{Q_{\ThPB}}(\t) \tag{Def. $Q_\eqsyn$}\\
    &\Longleftrightarrow \s \dot{\sim} \t \tag{Proposition \ref{prop:twoequivalence}}
  \end{align}
\end{proof}


\begin{corollary}\label{cor:finale}
  For all $c,d \in \Diag{\SigPRB}$, $\osem{c}=\osem{d}$ iff $\encoding{c} \eqsynbis \encoding{d}$.
\end{corollary}
\begin{proof}
  \begin{align}
    \osem{c}=\osem{d} &\Longleftrightarrow  \dsem{\encoding{c}}=\dsem{\encoding{d}} \tag{Proposition \ref{prop:encoding}}\\
    &\Longleftrightarrow \encoding{c} \eqsynbis \encoding{d}. \tag{Proposition \ref{prop:soundnesstapes}, Corollary \ref{cor:completenessPBPtapes}}
  \end{align}
\end{proof}